\section{Circuit and frame measurement}
\label{app:spice_readout}

\subsection{Circuit parameters}

The calculation uses BSIM3 transistor models for 180 nm CMOS, evaluated with
LTspice 17.2.4. Figure~\ref{fig:spice_circuit} gives the connections and gate
biases. Table~\ref{tab:spice_devices} gives the dimensions, including the
multiplicity $m$ of parallel devices. The pad capacitance is an external
load; intrinsic device capacitances are supplied by the transistor models.

The gate voltages of M3 and M4 are quoted as
\SI{\SpiceMthreeGateV}{V} and \SI{\SpiceMfourGateV}{V}, respectively.
Repeating the calculation with these rounded values changes the gain by
\SpiceBiasRoundingGainPercent\% and the ENC by
\SpiceBiasRoundingEncPercent\%.

The transient voltage is measured relative to a matched zero-input run.
The physical current has the square-pad boundary conditions specified in
Appendix~\ref{app:single_ion}; its integral is preserved in the tabulated
SPICE source. The gain $g$ is evaluated with the uniform arrival current,
the averaging windows and the frame difference defined in the main text.

\subsection{Contributions from neighbouring frames}

For equal-sized ion packets, the voltage left by a packet in the preceding
frame contributes \SpicePreviousPercent\% of the selected-frame signal
after the discharge correction. A packet collected in the following frame
contributes \SpiceFollowingPercent\%. Its ions induce current before reaching
the pad, and the voltage averaging window extends across the frame boundary.
Both contributions are calculated with the same averaging and discharge
correction used to recover the charge.

\subsection{Averaging and noise covariance}
\label{app:spice_noise}

Let $C_{\overline V}(u)$ be the covariance of averaged voltage readings
separated by $u$. For stationary noise with one-sided output spectrum $S_V$,
\begin{equation}
 C_{\overline V}(u)=\int_{f_\ell}^{f_h}S_V(f)
 \left|\frac{\sin(\pi fT_a)}{\pi fT_a}\right|^2
 \cos(2\pi fu)\,\dd f.
 \label{eq:spice_covariance}
\end{equation}
The same averaging window is applied to the transient signal. Squaring
Eq.~\eqref{eq:spice_noise_difference} and taking the average gives
\begin{equation}
 \sigma_D^2=(1+\alpha^2)C_{\overline V}(0)
           -2\alpha C_{\overline V}(\Delta t).
 \label{eq:spice_covariance_difference}
\end{equation}
Substituting Eq.~\eqref{eq:spice_covariance} yields
Eq.~\eqref{eq:spice_noise_integral}, since
\begin{equation}
 |1-\alpha e^{-2\pi\mathrm{i}f\Delta t}|^2
 =1+\alpha^2-2\alpha\cos(2\pi f\Delta t).
\end{equation}

The corrected difference retains a response $1-\alpha$ at zero frequency.
The lower integration limit $f_\ell$ therefore remains part of the noise
specification. The calculation uses the continuous average of the output
waveform and the circuit noise at its DC operating point.
ADC noise and quantisation must be included when a digitisation scheme is chosen.

The noise integral is checked by increasing the spectral sampling density
and the quadrature order. The signal is checked by reducing the transient
time step. Direct simulations of the adjacent integer ion counts bracket
the S/N threshold quoted in the main text.

Figure~\ref{fig:spice_averaging} shows how the averaging time affects the ENC.
Increasing the average from \SI{180}{\micro\second} to
\SI{350}{\micro\second} reduces the ENC from \SpiceScanFirstEnc\ to
\SpiceScanLastEnc\ electrons RMS. The signal per ion is recalculated for
each averaging time.

\begin{figure}[!htbp]
 \centering
 \includegraphics[width=0.88\linewidth]{spice_averaging.pdf}
 \caption{ENC versus averaging time for uniform ion arrivals over
 \SI{\SpiceFrameMs}{ms}, with an external pad capacitance of
 \SI{\SpicePadPf}{pF}. The highlighted point is the
 \SI{\SpiceAverageUs}{\micro\second} average used in this study.}
 \label{fig:spice_averaging}
\end{figure}
